% ============ Section II-A: Related work -- Fraud detection (Member 1) ============
% Goes inside Section II (Related Work). Member 2 writes II-B (concept drift),
% Member 3 writes II-C (algorithmic fairness); the summary table and the final
% "research gap" paragraph are joint.
\section{Related Work}
\label{sec:related}

\subsection{Fraud Detection with Machine Learning}
\label{sec:related-fraud}

\textit{Surveys.} Financial fraud detection has been reviewed many times. Al-Hashedi and
Magalingam \cite{alhashedi2021financial} reviewed data mining studies from 2009 to 2019 and grouped
them by fraud type and technique; credit card fraud was the most studied type. A later systematic
review by Ali~\textit{et al.} \cite{ali2022financial} reached the same conclusion and found support
vector machines and neural networks to be the most common models. Hilal~\textit{et al.}
\cite{hilal2022financial} focused on anomaly detection and on semi-supervised and unsupervised
methods, which are useful when labels are scarce. All three reviews point out that public data
for fraud research is limited, which is why many papers end up using the same few datasets.

\textit{Supervised models on transaction data.} Most recent supervised work uses credit card
transactions, very often the public European cardholder dataset. Alarfaj~\textit{et al.}
\cite{alarfaj2022credit} compared classical ML models with several convolutional neural networks on
that dataset and reported 99.9\% accuracy and an F1-score of 85.71\% for their best deep model.
Taha and Malebary \cite{taha2020intelligent} tuned LightGBM with Bayesian optimisation on two public
card datasets and reported 98.40\% accuracy but an F1-score of only 56.95\%, which shows how far
apart accuracy and fraud-focused metrics can be on imbalanced data. Zhao~\textit{et al.}
\cite{zhao2024improved} changed the LightGBM loss to deal with extreme imbalance and combined it with
oversampling. Gradient-boosted trees appear again and again in this line of work, which agrees with
general benchmarks showing that tree ensembles still beat deep networks on typical tabular data
\cite{grinsztajn2022tree, shwartzziv2022tabular}. The weakness shared by these studies is the evaluation:
accuracy is given as a headline metric even though fraud is rare, and none of them reports how
performance changes from one time period to the next.

\textit{Realistic evaluation and feature design.} A smaller group of papers works with data from
real fraud detection systems and tries to copy how they operate. Dal Pozzolo~\textit{et al.}
\cite{dalpozzolo2018credit} described a fraud detection system with more than 75 million
transactions over three years and showed that class imbalance, concept drift and delayed labels all
have to be handled together; they also argued that alert precision, not accuracy, is what
investigators care about. Carcillo~\textit{et al.} \cite{carcillo2021combining} added unsupervised
outlier scores as extra features for a supervised model; some scores helped, but using many of
them at once did not. Correa
Bahnsen~\textit{et al.} \cite{bahnsen2016feature} showed that aggregated features, such as how many
transactions a customer made in the last hours or days, add a lot to a fraud model. The BAF data
contains the same kind of aggregated ``velocity'' features at the application level. These papers
are close to our setting, but their data is private, so their results cannot be reproduced.

\textit{Account opening fraud and the BAF suite.} Account opening fraud is studied much less than
transaction fraud, mainly because no large public dataset existed before BAF
\cite{jesus2022turning}. BAF was generated from a real bank's application data with a tabular
generative model and extra noise for privacy, and it keeps the class imbalance, the monthly changes
and the group differences of the original data. Its authors trained 100 LightGBM configurations with
random search, trained on months 0--5 and tested on months 6--7, and reported the true positive rate
at a 5\% false positive rate and the ratio of false positive rates between age groups. Sun~\textit{et
al.} \cite{sun2025objective} compared logistic regression, an SVM, random forest, LightGBM and a GRU
on the same Base dataset. At the 0.5 threshold the classical models caught only 1--6\% of fraud,
while a class-weighted GRU caught 78\% with about 5\% precision. Their main point, that the training
objective and the operating point matter as much as the model type, supports reporting results at a
fixed false positive rate. However, they used a stratified random split instead of the month
column, tuned the classical models on a 10{,}000-row sample, and did not look at fairness.

\textit{Gap.} Taken together, the fraud detection literature either uses realistic but private
data, or public data with random splits and accuracy-style metrics. On BAF, the only public
account-opening benchmark, no study we found compares several model families under one time-based
protocol in which the 5\% FPR threshold is chosen without looking at the test months. Our study
fills this part of the gap; Sections~\ref{sec:related-drift} and \ref{sec:related-fair} cover the
drift and fairness sides.
